\section{Executive Summary}

\begin{calloutbox}{Research Vision}
My research program lies at the convergence of \textbf{computational neuroscience, solid-state VLSI circuits, and emerging non-volatile memories}. Over the next five years, my goal is to design, manufacture, and deploy neuromorphic computing systems that achieve biological-level energy efficiency and real-time adaptability for edge computing applications. By co-designing hardware substrates with algorithms that emulate brain-like plastic changes, my research addresses the physical scaling limitations and computational bottlenecks of traditional Von Neumann architectures.
\end{calloutbox}

Modern artificial intelligence (AI) has achieved unparalleled success, but this has come at the expense of massive carbon and energy footprints. While biological brains operate on a mere 20 Watts of power, state-of-the-art silicon accelerators require hundreds of Watts to perform comparable tasks. This disparity stems from the separation of memory and processing units in classical computing architectures, which leads to substantial data movement energy costs. 

To bridge this gap, my research program establishes a multi-disciplinary co-design framework spanning three core pillars:
\begin{enumerate}
    \item \textbf{Device-Level In-Memory Computing:} Using memristive and phase-change materials to perform matrix-vector operations directly within storage arrays.
    \item \textbf{Algorithm-Level Event-Driven Execution:} Building spiking neural networks (SNNs) that process information asynchronously, mimicking biological action potentials.
    \item \textbf{System-Level Online Adaptability:} Implementing hardware-in-the-loop plasticity rules that enable silicon agents to learn and recover from failures in real-time.
\end{enumerate}

To date, my work has resulted in 12 publications in premier journals and conferences (including IEEE ISSCC, DAC, and NeurIPS), with direct validation of hardware designs in 28nm CMOS processes. My open-source neuromorphic simulator, \texttt{NexaSpike}, is currently utilized by more than 40 academic groups and industrial labs to evaluate edge computing architectures.
